\chapter{Objetivos}
\label{cap:objetivos}

\section{Objetivo geral}

Analisar comparativamente as especificações técnicas CNC-OMBR-MAT-18-0126-EDCE e CNC-NDBR-DBR-25-1580-EDBR, identificando as alterações nos requisitos de fornecimento de energia elétrica aplicáveis a empreendimentos com múltiplas unidades consumidoras na área de concessão da Enel Distribuição Ceará e avaliando seus possíveis impactos técnicos sobre o projeto elétrico do Edifício Palatium, localizado em Sobral, Ceará.

\section{Objetivos específicos}

\begin{itemize}
  \item Identificar os requisitos técnicos estabelecidos pelas especificações 0126 e 1580 aplicáveis à área de concessão da Enel Distribuição Ceará;
  \item Comparar e classificar as diferenças entre as especificações, considerando os critérios de atendimento, dimensionamento, proteção, medição, aterramento, documentação e demais requisitos técnicos de conexão coletiva;
  \item Sistematizar as diferenças normativas identificadas por meio de uma matriz comparativa;
  \item Identificar, entre as alterações normativas, aquelas aplicáveis às características e às soluções técnicas adotadas no projeto elétrico do Edifício Palatium;
  \item Avaliar os possíveis impactos técnicos, construtivos e documentais decorrentes da aplicação dos requisitos da especificação 1580 ao projeto elétrico do empreendimento, considerando as informações disponíveis e as limitações da análise.
\end{itemize}

\section{Organização do trabalho}

O presente trabalho está estruturado em sete capítulos.

O Capítulo 1 apresenta a contextualização do tema, a problemática da pesquisa e a justificativa para seu desenvolvimento.

O Capítulo 2 estabelece o objetivo geral e os objetivos específicos, além de apresentar a organização do trabalho.

O Capítulo 3 apresenta os trabalhos relacionados ao tema da pesquisa, considerando estudos sobre instalações elétricas, critérios de fornecimento de energia elétrica e análise de requisitos normativos.

O Capítulo 4 apresenta a fundamentação teórica necessária à compreensão dos aspectos técnicos e regulatórios relacionados ao fornecimento de energia elétrica em empreendimentos com múltiplas unidades consumidoras.

O Capítulo 5 descreve a metodologia empregada, o conjunto documental utilizado, a caracterização do empreendimento e os procedimentos de comparação normativa e avaliação dos impactos técnicos.

O Capítulo 6 apresenta e discute os resultados da comparação entre as especificações técnicas e da aplicação das alterações normativas ao estudo de caso.

Por fim, o Capítulo 7 apresenta as conclusões da pesquisa, suas contribuições, limitações e sugestões para trabalhos futuros.
